\section{The retractive example and remaining questions}\label{sec:finish}

We now turn to the shear construction of Theorem~\ref{thm:shear}. The base group is finitely
presented and diffuse, its normalized line entries form free bases,
and its line matrices satisfy a strict rank bound. These are exactly
the hypotheses of the attachment theorem.

\begin{proof}[Proof of Theorem~\ref{thm:shear}]
Fix a prime \(p\) and a power \(Q=p^a>Q_0\), with \(Q_0\) as
in~\eqref{eq:qbound}, and construct the finite complex \(X\) above
from the incidence data over \(\F_Q\). Write \(G_p=\pi_1X\).
By Corollary~\ref{cor:base-diffuse}, \(G_p\) is diffuse and hence has
unique product. Lemma~\ref{lem:free-core} supplies free subgroups
\(F_i\), and the normalization following that lemma puts the
matrices \(Y_i\) in the form required by Theorem~\ref{thm:principle}.
Proposition~\ref{prop:rank-budget} gives
\[
 \sum_i\rk_D\phi(Y_i)<Nr/4<Nr
\]
for every homomorphism \(\F_p[G_p]\to D\) to a division ring.

Apply Theorem~\ref{thm:principle} with \(K=\F_p\). The resulting
group
\[
 \Lambda_p=G_p*_{F_i,\ i\in I}(\Z^2\rtimes F_i)
\]
is finitely presented and diffuse, and \(\F_p[\Lambda_p]\) embeds in
no division ring. Unique product makes \(K[\Lambda_p]\) a domain for
every field \(K\). If \(\operatorname{char}K=p\), any injection of
\(K[\Lambda_p]\) into a division ring would restrict to an injection
of \(\F_p[\Lambda_p]\). This is impossible.

Lemma~\ref{lem:kernel} identifies the retraction kernel as a free
product of copies of \(\Z^2\). It is bi-orderable by the free-product
ordering theorem~\cite{Vinogradov}. A subgroup of a left-orderable
group is left-orderable, so left-orderability of \(\Lambda_p\) implies
that of \(G_p\). Conversely, an extension of left-orderable groups is
left-orderable: order an element by its image in the quotient, and,
when that image is the identity, by a chosen order in the kernel.
The resulting positive cone is closed under multiplication and
partitions the nonidentity elements with its inverse. Hence a left
order on \(G_p\), together with one on the kernel, gives a left order
on \(\Lambda_p\). This proves the remaining assertions.
\end{proof}

\subsection{A finite presentation without unspecified choices}

The proof specifies a presentation, rather than just an existence
class. Fix a prime \(p\). Let \(a\) be the least exponent with
\(Q=p^a>Q_0\), and let \(\pi\) be the first monic irreducible
polynomial of degree \(a\) over \(\F_p\) in the lexicographic order
of its coefficient residues. This fixes \(\F_Q\), its numbering, and
\(S\). The lexicographic rule after
Lemma~\ref{lem:incidence} determines the marked lines. The integer
formula in Section~\ref{sec:lines} determines the two blocks at each
regular point, and~\eqref{eq:coordinates} determines their labels.

Order the vertices and edges of \(X\) by their displayed indices,
and choose the first spanning tree \(T\) in this order. Its paths
give the paths \(p_T\) used to define the based edge elements.
A presentation of \(G_p\) has one generator \(g_e\) for each oriented
edge outside the spanning tree, with \(g_e=1\) on tree edges,
\(g_{\bar e}=g_e^{-1}\), and the triangle relations
\[
 g_{L_{b,u,\alpha_b(i)+v}}\,g_{x_{iuv}}
 =g_{R_{b,v,\beta_b(i)+u}}.
\]
For each \(i\in I\), form the words \(e_{i,uv}\) in
\eqref{eq:free-generators}. Add two generators \(s_i,\gamma_i\) and
the relations
\begin{equation}\label{eq:final-presentation}
 \begin{aligned}
 [s_i,\gamma_i]&=1,\\
 [s_i,e_{i,uv}]&=1,\\
 \gamma_i e_{i,uv}\gamma_i^{-1}
 &=s_i^{\delta_{uv}}e_{i,uv}
 \qquad(1\le u,v<r).
 \end{aligned}
\end{equation}
This is a finite presentation of \(\Lambda_p\).
The search for the incidence data is finite and terminates by
Lemma~\ref{lem:incidence}; no claim of computational feasibility is
needed for this specification.

The base \(G_p\) is word-hyperbolic. The final group \(\Lambda_p\)
contains each attachment's copy of \(\Z^2\), so it is not
word-hyperbolic. Its useful retained structure is the retraction
\[
 1\longrightarrow
 *_{i,\,gF_i\in G_p/F_i}g\Z^2g^{-1}
 \longrightarrow\Lambda_p\longrightarrow G_p\longrightarrow1.
\]
It is this structure, rather than curvature of the enlarged group,
that preserves diffuseness.

\subsection{What the attachments add}

The base already excludes a natural class of division-ring embeddings.
We record this to distinguish the rank condition supplied by an
embedding from the one forced by the new group relations.
For OpenAI's group the corresponding statement is
\cite[Theorem~1.4 and Corollary~6.5]{DM}.

For an embedding \(K[H]\to D\), write \(D_C\) for the division
closure of \(K[C]\) when \(C\le H\).
A \emph{Linnell division ring} is a division ring generated by
\(K[H]\) in which, for every subgroup \(C\le H\), elements in
distinct right cosets of \(C\) are left linearly independent over
\(D_C\); see~\cite[Section~2.3]{JZL} and~\cite[Definition~2.2]{FSP}.

\begin{proposition}\label{prop:base-linnell}
For every field \(K\) of characteristic \(p\), the domain \(K[G_p]\)
has no Linnell division ring.
\end{proposition}

\begin{proof}
Suppose such a division ring exists. Fix a line, omit its label, and
put \(M_k=(e_{uv}-1)_{1\le u,v\le k}\).
We prove that each \(M_k\) is invertible. For \(k=1\), this follows
from injectivity and \(e_{11}\ne1\). For the inductive step, partition
\[
 M_k=\begin{pmatrix}M_{k-1}&b\\c&e_{kk}-1\end{pmatrix}.
\]
Let \(C\) be the subgroup generated by all the free basis elements
of this line except \(e_{kk}\). Gaussian elimination over \(D_C\)
shows that the entries of \(M_{k-1}^{-1}\), as well as those of
\(b\) and \(c\), lie in \(D_C\). The Schur complement is
\[
 e_{kk}+d,\qquad d=-1-cM_{k-1}^{-1}b\in D_C.
\]
It is nonzero: \(e_{kk}\notin C\), so \(1\) and \(e_{kk}\) belong
to distinct right cosets and are independent over \(D_C\).
Block elimination makes \(M_k\) invertible. The normalization of
\(Y_i\) then gives \(\rk_D Y_i=r\) for every \(i\), contradicting
Proposition~\ref{prop:rank-budget}.
\end{proof}

For an arbitrary division-ring embedding of the base algebra, that
coset independence is not available. The attachments replace it by
the conjugation calculation of Lemma~\ref{lem:rank-forcing}.
Thus Theorem~\ref{thm:shear} does not infer ordinary nonembeddability
from the failure of a Linnell embedding.

Whether \(\F_p[G_p]\) itself embeds in some division ring remains
open here. An embedding would give a domain that embeds in a division
ring but admits no Linnell division ring. A proof of nonembeddability
would give a smaller hyperbolic example using the base itself, without
the permutation graph of Section~\ref{sec:hyperbolic}.
Neither construction decides whether \(\Z[\Lambda_p]\) or
\(\mathbb{Q}[\Lambda_p]\) embeds in a division ring, and we likewise
make no such decision for \(\Z[\Gamma_p]\) or
\(\mathbb Q[\Gamma_p]\); the integral and
characteristic-zero versions of Kourovka Problem~1.5 remain open.

The kernel of the retraction \(\Lambda_p\to G_p\) is a free product of
copies of \(\Z^2\), hence bi-orderable~\cite{Vinogradov}. Thus
\(\Lambda_p\) is left-orderable if and only if \(G_p\) is. We do not know
whether \(G_p\) is left-orderable. If it is, \(\F_p[\Lambda_p]\) answers
Mal'cev's problem (Kourovka Problem~1.6~\cite{Kourovka}) negatively in
characteristic~\(p\).

The hyperbolic construction answers the positive-characteristic field
version of Kourovka Problem~1.5 with word-hyperbolic groups. Relative
to the shear construction, its additional ingredients are the
permutation average and the concentration proof for the line subgroups.
It does not settle the orderability of the base or Mal'cev's problem;
the retractive example retains the precise reduction of its
left-orderability to that of \(G_p\).

\paragraph{Changes in version 1.1.0.}
This version adds the Lean verification of the nonembedding clauses of
Theorems~\ref{thm:main} and~\ref{thm:shear}, with an appendix describing
its scope and checks. The mathematical statements and their proofs in
the text are unchanged.

\section*{How this work was done}

The authors directed the research. The
OpenAI models GPT-6 Astra and GPT-6.1 Sol, directed by Nidhal Mghirbi,
assisted with construction, proof exploration, checking and writing,
including the shear attachment and base geometric arguments. Claude
Opus~5.5 (Anthropic) checked those underlying proofs refutation-first,
and a semi-blind GPT-6.1 Sol replication proved their main steps again.
The extension from sufficiently large primes to every prime was found
in a GPT-6.1 Sol research run and checked by Claude Opus~5.5.
The hyperbolic construction was found in a GPT-6.1 Sol research run
and proved again in an independent blind GPT-6.1 Sol replication, which
supplied the concentration proof. Claude Opus~5.5 then checked the
hyperbolic construction refutation-first, re-deriving each load-bearing
step before reading either run's argument; its repairs, including the
full write-up of the puncture-lemma proof and the combination-theorem
citation, are incorporated here.

The Lean development was written by GPT-6.1 Sol agents in Codesk
sessions, directed by the authors. An independent blind fidelity read
found the statements faithful to the paper; they were then frozen,
before the proofs were written. The recorded package passed the default
build, the full axiom audit, a sandboxed comparator against a
Mathlib-only challenge, and proof-closure replay by \texttt{leanchecker}
and \texttt{nanoda}, with a negative control omitting
\texttt{Classical.choice}. Release lint and a fresh-clone cache retrieval,
build and audit also passed, and the shipped folder passed them again
from a clean copy. Appendix~\ref{app:lean} gives the scope, pins and
check results.

The authors take responsibility for
the statements and proofs. The proofs are given in full; the internal
checks are not a substitute for them.

\section*{Acknowledgements}
We thank Andrei Jaikin-Zapirain for helpful correspondence, and for his
questions, which motivated this work.
